\begin{frame}{Research Gaps}{Limitations in Evaluating Interpretability}
  \begin{itemize}
    \item \textbf{Reconstruction vs. Semantic Grounding}: High SAE sparsity and reconstruction fidelity do not guarantee that the nearest text label correctly describes the feature[cite: 3].
    \item \textbf{Dependence on External Resources}: The discovery pipeline relies heavily on external tools (ontology vocabularies, generative phrasing, LLM judges)[cite: 3]. Biases in these components can be mistaken for properties of the learned representation[cite: 3].
    \item \textbf{Selective Evidence in Reports}: Radiology reports provide partial evidence[cite: 3]. Unmentioned concepts are not necessarily absent, causing external judges to confuse omission with contradiction[cite: 3].
    \item \textbf{Independent Labeling}: Assigning labels independently omits causal and anatomical relationships between findings[cite: 3]. We explore semantic clustering to identify useful macro-groups[cite: 3].
  \end{itemize}
  
  \note{
    We identified four major research gaps that guide our analysis. 
    First, mathematical reconstruction fidelity does not ensure semantic correctness; the geometry of the embedding space heavily constrains the naming process[cite: 3]. 
    Second, the entire pipeline is dependent on external resources like UMLS ontologies and generative models[cite: 3]. Errors generated by these tools can easily be misattributed to the SAE itself[cite: 3]. 
    Third, radiology reports are selective by nature[cite: 3]. Radiologists often omit normal anatomy, which an automated judge might wrongly penalize as a contradiction[cite: 3]. 
    Lastly, treating concepts in isolation ignores their real-world anatomical relationships, a gap we attempt to address later through hierarchical clustering[cite: 3].
  }
\end{frame}